\chapter{What Are Cellular Automata?}

\section{Introduction}

Back in 1971 an article appeared in the Mathematical Games section of
\emph{Scientific American} describing the Game of Life, created by
Cambridge mathematician John Conway.\footnote{Martin Gardner's column
actually introduced Life in October 1970; follow-up columns appeared
through 1971.}  Conway had been looking for a simple discrete dynamical
system capable of self-replication.  Soon after the article appeared,
computers the world over began devoting unprecedented amounts of time
to exploring Life's ramifications, often to the chagrin of computing
facility administrators.

Conway's Game of Life is one of an infinite class of mathematical
systems known as \emph{cellular automata} (CA).  Creating a CA is
creating a simple universe, with its own space-time structure and its
own laws of physics.  One sets up the initial state (perhaps using
randomization), defines the laws of physics, starts time's progression,
and beholds the subsequent evolution of the universe.  When playing God
in this way, one is often surprised and delighted by what ensues.

CA were first introduced by John von Neumann and Stanis\l{}aw Ulam in
the 1940s, who called them \emph{cellular spaces}; they had been
looking for simple mathematical models of biological systems,
self-reproduction in particular.  CA have since served as mathematical
models of phenomena from remarkably diverse disciplines.  Any system
analyzable in terms of large numbers of discrete elements with local
interactions is amenable to being modeled as a CA, and in such systems
simple interaction rules between components tend to give rise to very
complex \emph{emergent} behavior.

CA can yield discrete approximations to the solutions of systems of
differential equations, in terms of which much of the macroscopic
physics of our world is expressed.  CA models have been applied to
fluid dynamics, plasma physics, chemical systems, the growth of
dendritic crystals, economics, two-directional traffic flow, image
processing and pattern recognition, parallel processing, random number
generation, and even the evolution of spiral galaxies.  Biological
applications include the functioning of groups of cells, heart
fibrillation, neural networks, and ecosystems.

The module \texttt{cadyn.ecosystem} uses a CA to simulate a
three-species ecosystem, capturing the dynamics typical of such
systems.  Mathematicians have traditionally used generalizations of the
Volterra--Lotka differential equations to model the population dynamics
of predator/prey systems.  CA provide a more intuitive model ---
indeed, a CA models the ecosystem itself, not merely the population
counts of the species.

CA also provide a useful mathematical model of massively parallel
multiprocessor systems.  Each cell may be considered a processor, its
finitely many states the states of the processor, and its neighborhood
the set of directly connected processors.  (Replace ``processor'' with
``neuron'' and the same description covers a neural network.)  How to
get such a system to perform useful computation, making optimal use of
all that parallel computing power, remains a central problem of
computer science.  CA experiments have provided much-needed insight
into how simple local interactive dynamics give rise to complex
emergent global behavior.

Similar questions --- simple local rules yielding complex emergent
global behavior --- are central to many fields.  Some important
unanswered questions of this type:

\begin{itemize}
\item How did life emerge from the interactions of chemicals in the
  pre-biotic conditions of the primordial earth?
\item How does morphogenesis proceed in a developing embryo?  How do
  the relatively simple, chemically modulated interactions of
  undifferentiated cells result in a globally orchestrated complex of
  specialized organs and tissues?
\item How does consciousness emerge from the relatively simple
  interactions of billions of neurons?
\item How does cultural evolution unfold from the interaction of its
  underlying milieu of concepts and beliefs?
\item How do relatively simple economic behavior patterns of
  individuals give rise to the dynamics of a whole economy?
\end{itemize}

CA computer experiments and simulations provide insight into such
problems, and a conceptual framework in which theoretical progress may
be made on the phenomenon of emergence in general.  The Santa Fe
Institute was formed to address just these types of problems from a
variety of fields, and CA remain a major focus of research on complex
systems.

As single-processor computer systems yield to massively parallel ones,
simulation of CA becomes ever easier, and practical applications
abound.  The inherently parallel structure of CA makes them
particularly suited to implementation on multiprocessor hardware.  The
MIT Physics of Computation group developed a cellular automata machine
(CAM) specially designed to run CA programs, successfully applied to
problems in physics and in medical imaging.

\expand{The trajectory sketched above has largely come to pass.  The
GPUs that now power scientific computing (and machine learning) are
precisely the massively parallel, locally connected hardware for which
CA are natively suited; a modest laptop today evolves in seconds the
$500\times500$, 250-generation experiments that in 1997 required an
overnight run.  The lattice-gas methods mentioned above matured into
the lattice Boltzmann methods now standard in computational fluid
dynamics.  Toffoli and Margolus's CAM machines were retired, but their
ideas about reversible, physics-like computation live on in research on
energy-efficient and quantum computing.  And in 2002 Stephen Wolfram's
\emph{A New Kind of Science} brought CA to a wide audience; see
Chapter~\ref{ch:classification}.  The ecosystem strand, meanwhile, grew
a second act: Chapter~\ref{ch:evolution} adds \emph{movement} and
competing \emph{strategies} to the predator/prey CA and turns it on a
question from the science of mind --- whether evolution favors accurate
perception.}

\section{Basic Definitions}

Cellular automata are a class of discrete dynamical systems consisting
of an array of cells of some dimension $n$.  Each cell is in one of $k$
distinct states at each tick of a global clock, and at each tick every
cell may change state as determined by the \emph{transition rules} of
the particular CA.  The transition rules prescribe how a cell changes
state as a function of its current state and the states of its
\emph{neighbors}; which cells constitute the neighborhood of a given
cell must be specified explicitly.

More precisely, a CA consists of:

\begin{enumerate}
\item \textbf{A lattice of cells}, each of which is in one of a finite
  number of states at each discrete time.  The lattice may be of any
  dimension and is in principle of infinite extent: a 1-dimensional CA
  has a cell at each integer point of the real line, a 2-dimensional CA
  a cell at each point of the plane with integer coordinates, and so
  on.  In practice the infinite lattice is made effectively finite,
  usually by imposing a \emph{wrap-around} (periodic) condition: in one
  dimension the line becomes a circle, in two dimensions the plane
  becomes a torus.
\item \textbf{A neighborhood} defined for each cell.  For
  \emph{homogeneous} CA the same relative neighborhood is used for
  every cell.
\item \textbf{Transition rules} specifying the dynamics: each cell
  passes from its current state to its next state based on its current
  state and those of its neighbors, all cells transforming
  synchronously at the discrete ticks of the clock.
\end{enumerate}

Let $n$ be the dimension of the lattice, $k$ the number of states, $T$
the transition-rule function, $C_t(i_1,\dots,i_n)$ the state of the
cell at position $(i_1,\dots,i_n)$ at time $t$, and
$N_t(i_1,\dots,i_n)$ the tuple of values (in a fixed order) of the
cells neighboring that position at time $t$.  Then the dynamics of the
CA are completely specified by the initial states $C_0$ of all cells
together with the recursion
\begin{equation}
  C_{t+1}(i_1,\dots,i_n) \;=\; T\bigl(N_t(i_1,\dots,i_n)\bigr).
  \label{eq:ca-recursion}
\end{equation}

As is usual with formal definitions, this is best understood through
simple examples.  The next chapter examines the case in which the
lattice has dimension zero --- there is only one cell!  We shall see
that even this extremely simple case holds some surprises.
